
\section{Introduction}
\label{sec:introduction}

\noindent
Treemaps provide mappings of tree-structured, hierarchical data by means of space-constrained, recursively nested sets of rectangles that express the parent-child relationship among nodes and whose sizes are proportional to some computed per-node weights \cite{Shneiderman2009}.
Data associated with nodes, here referred to as \textit{attributes}, can be mapped by means of the 
\textit{visual variables} of treemaps, \eg size, color, texture, and shading \cite{Carpendale2003}.

\textit{2.5D treemaps} extend treemaps by using the third dimension to extrude rectangles to \textit{blocks} while still having the regular treemap's two-dimensional reference space and layout. 
Technically, they are represented as interactive 3D models and often ambiguously denoted by \textit{3D treemaps}. 
On the one hand, blocks facilitate emphasizing the hierarchical nesting of inner nodes and offer size and color as principal visual variables. 
On the other hand, blocks also provide another strong visual variable, namely height, 
which, in consequence, allows us to independently map three attributes (size, color, height) 
and to visually relate these attributes in a single 2.5D treemap view. 

Many 2D and 2.5D treemap variants are applied in a growing number of applications and systems to visualize multivariate, large tree-structured data of, \eg 
file systems \cite{Shneiderman1992}, 
software systems \cite{Wettel2008}, 
business data \cite{Vliegen2006}, 
stock markets \cite{Wattenberg1999}, or 
controller performance data \cite{Mitchell2004}.
For example, visual software analytics tools build software maps based on 2.5D treemaps 
to provide insight into large software artifact hierarchies and their characteristics 
by relating code modules (nodes) with associated data about 
code metrics, software quality, development activity as well as budget and resource planning \cite{Bohnet2011}.

As Wood et al. \cite{Wood2012} point out, \textit{sketchiness} 
\enquote{offers different ways of communicating ideas of narrative, purpose, ownership, accuracy and aesthetic} 
and it \enquote{may be reliably used as a visual variable on an ordinal scale, but that caution should be exercised when representing interval or ratio scale data}.
Uncertainty, imprecision, or vagueness data, 
mostly represented by \textit{ordinal variables with small range}, 
can be found in many treemap application scenarios.
Of course, this kind of attribute data can be mapped to common visual variables. 
However, if we depict this data by means of sketchiness, 
which provides a fourth `strong' visual variable for 2.5D treemaps, 
we thereby retain the other visual variables.
Sketchiness also allows us to qualify attribute mappings 
by mapping uncertainty in a generally understandable way. 
For example, in software maps we could superimpose the attribute `test coverage' (sketchiness) to the mappings of lines-of-code (size), software complexity (height) and frequency of changes (color) to emphasize those modules that are difficult to comprehend, have been significantly altered recently \textit{and} that are rarely tested. 
This paper provides the following contributions:
\begin{enumerate}[leftmargin=16pt]
	\item We discuss a design space of sketchiness for 2.5D treemaps \wrt rendering of sketchy outlines and surfaces, which is explained in detail for three distinct styles: Biro, (colored) Pencil, and Marker. %
	\item We evaluate sketchiness in terms of a visual 2.5D treemap variable and discuss its selective, associative, quantitative, as well as order and length characteristics.
\end{enumerate}

